\section{From energy recovery to path tightness}\label{sec:tightness}

Finite-dimensional convergence does not control excursions between observation
times. We obtain that control from the local energy recovery in
Proposition~\ref{prop:inputs}. First we extend buffered cutoff recoveries to
the whole graph. A reversible maximal inequality then makes their resolvent
approximations accurate along the entire stationary path. The bounded
generators of those approximations will control motion after stopping times.

\subsection{Uniform recovery and a stationary maximal inequality}

\begin{lemma}[Uniform cutoff recovery]\label{lem:uniform-recovery}
Let $\chi:\Sph\to[0,1]$ be smooth, and put $f=1-\chi$. Suppose
$\operatorname{supp}\chi\subset V\Subset W\Subset U$ for buffered conformal
chart regions to which the local recovery statement of
Proposition~\ref{prop:inputs} applies. There are real vertex functions $w_n$
such that
\begin{equation}\label{eq:uniform-recovery}
 \norm{w_n-f\circ z_n}_{\infty,V_n}\longrightarrow0,
 \qquad \limsup_n\E_n(w_n)\le\E_h(f).
\end{equation}
\end{lemma}

\begin{proof}
Apply local recovery to $\chi$ on $U$, obtaining $\widetilde\chi_n$ uniformly
close to $\chi\circ z_n$ at vertices in $W$. Choose $e_n\downarrow0$
bounding these uniform errors; decreasingness can be achieved by taking tail
suprema. Define
\[
 \chi_n(v)=
 \begin{cases}
 \min\{1,\max\{0,\widetilde\chi_n(v)-e_n\}\},&z_n(v)\in W,\\
 0,&z_n(v)\notin W,
 \end{cases}
 \qquad w_n=1-\chi_n.
\]
The clipping map is $1$-Lipschitz, and $\chi$ vanishes outside $W$, so
$\norm{\chi_n-\chi\circ z_n}_\infty\le2e_n$. More importantly,
$\chi_n(v)>0$ implies $\chi(z_n(v))>0$. Thus an edge with a nonzero difference
of $\chi_n$ has an endpoint in $\operatorname{supp}\chi\Subset V$.
The vanishing edge mesh and the buffer ensure that every such edge is
eventually wholly in $W$, with both endpoints there. Its difference is then
bounded in absolute value by the corresponding difference of
$\widetilde\chi_n$. Local recovery therefore gives
\[
 \limsup_n\E_n(w_n)
 =\limsup_n\frac12\sum_e(d_e\chi_n)^2
 \le\frac12\int_U\abs{\nabla\chi}^2\dd\mathrm{vol}
 =\E_h(f).
\]
The last equality uses the support condition and conformal invariance of
the two-dimensional Dirichlet integral. This proves
\eqref{eq:uniform-recovery}, including the factor $1/2$.
\end{proof}

The following estimate is a finite-state analogue of the forward--backward
martingale argument of Lyons and Zheng \cite{LyonsZheng1988}.
We include its proof to handle jumps explicitly and to make the uniformity
in the chain clear.

\begin{lemma}[Stationary maximal inequality]\label{lem:stationary-maximal}
Let $Y$ be a continuous-time Markov chain on a finite state space, with
generator $-A$ and reversible probability measure $\pi$, started from $\pi$.
For every real function $v$ on that space and every $T>0$,
\begin{equation}\label{eq:stationary-maximal}
 \EE\sup_{0\le t\le T}\abs{v(Y_t)}^2
 \le 2\norm{v}_{L^2(\pi)}^2+40T\langle v,Av\rangle_{L^2(\pi)}.
\end{equation}
The constants do not depend on the state space or its jump rates.
\end{lemma}

\begin{proof}
Write $q(x,y)$ for the off-diagonal jump rates. The process
\[
 N_t=v(Y_t)-v(Y_0)+\int_0^t Av(Y_s)\dd s
\]
is a square-integrable martingale starting from zero. Its quadratic variation
and stationarity give
\begin{align*}
 \EE\abs{N_T}^2
 &=T\sum_x\pi(x)\sum_{y\ne x}q(x,y)(v(y)-v(x))^2\\
 &=2T\langle v,Av\rangle_{L^2(\pi)}.
\end{align*}
Over $[0,T]$, reverse the path using its right-continuous modification in
reversed time. Reversibility and stationarity give the same chain law, so its
corresponding martingale $N'$ satisfies the same second-moment identity.
At every fixed $t\in[0,T]$, almost surely there is no jump at $t$ or at either
deterministic endpoint. Subtracting the forward and backward formulas yields
\[
 2\bigl(v(Y_t)-v(Y_0)\bigr)=N_t-(N'_T-N'_{T-t}).
\]
Use this identity first at a countable dense set of times, together with
$T$, and then use right continuity. It follows that
\[
 \sup_{t\le T}\abs{v(Y_t)}
 \le \abs{v(Y_0)}+\frac12\sup_{t\le T}\abs{N_t}
                       +\sup_{t\le T}\abs{N'_t}.
\]
If the two martingale suprema are $a$ and $b$, the square of the right side is
at most $2\abs{v(Y_0)}^2+a^2+4b^2$. Doob's $L^2$ inequality bounds each of
$\EE a^2$ and $\EE b^2$ by
$8T\langle v,Av\rangle_{L^2(\pi)}$, proving
\eqref{eq:stationary-maximal}.
\end{proof}

\subsection{Resolvent approximation along the whole path}

\begin{lemma}\label{lem:resolvent-approximation}
For a cutoff $f=1-\chi$ as in Lemma~\ref{lem:uniform-recovery}, write
$g_n=f\circ z_n$ and define $u_{n,\alpha}$ and $u_\alpha$ by
\eqref{eq:resolvent-functions}. For every $T>0$,
\begin{equation}\label{eq:pathwise-resolvent-approximation}
 \lim_{\alpha\to\infty}\limsup_n
 \EE\sup_{0\le t\le T}
       \abs{u_{n,\alpha}(Y^n_t)-g_n(Y^n_t)}=0.
\end{equation}
For every $n$ and $\alpha>0$, moreover,
\begin{equation}\label{eq:bounded-generator}
 0\le u_{n,\alpha}\le1,
 \qquad \norm{G_nu_{n,\alpha}}_\infty\le\alpha.
\end{equation}
\end{lemma}

\begin{proof}
Take the recoveries $w_n$ from Lemma~\ref{lem:uniform-recovery}. Their uniform
approximation, together with the tour convergence, gives $J_nw_n\to Jf$ in
$H$. For fixed $\alpha$, Equation~\eqref{eq:resolvent-graph-convergence}
therefore implies
\begin{align}
 &\limsup_n\bigl(
   \norm{w_n-u_{n,\alpha}}_{L^2(\mu_n)}^2
       +\E_n(w_n-u_{n,\alpha})\bigr)\notag\\
 &\hspace{15mm}\le
   \norm{f-u_\alpha}_{L^2(\mu_h)}^2+\E_h(f-u_\alpha).
 \label{eq:form-norm-approximation}
\end{align}
To see the energy assertion explicitly, expand
\[
 \E_n(w_n-u_{n,\alpha})
 =\E_n(w_n)-2\langle w_n,G_nu_{n,\alpha}\rangle_{L^2(\mu_n)}
            +\langle u_{n,\alpha},G_nu_{n,\alpha}\rangle_{L^2(\mu_n)}.
\]
The last two terms converge by \eqref{eq:resolvent-graph-convergence}, and
the first has limsup at most $\E_h(f)$. Since $u_\alpha$ lies in the operator
domain of $G$, the resulting expression is the form expansion of
$\E_h(f-u_\alpha)$.

The right side of \eqref{eq:form-norm-approximation} tends to zero. Indeed,
if $\nu_f$ is the spectral measure of $G$ at $f$, then it equals
\[
 \int_{[0,\infty)}(1+\lambda)
       \left(\frac{\lambda}{\alpha+\lambda}\right)^2\dd\nu_f(\lambda).
\]
Here $f$ belongs to the form domain, so $1+\lambda$ is integrable against
$\nu_f$; dominated convergence applies. Lemma~\ref{lem:stationary-maximal},
applied with $A=G_n$ to $w_n-u_{n,\alpha}$, now gives
\[
 \lim_{\alpha\to\infty}\limsup_n
 \EE\sup_{t\le T}\abs{w_n(Y^n_t)-u_{n,\alpha}(Y^n_t)}^2=0.
\]
Cauchy--Schwarz and the uniform bound on $w_n-g_n$ yield
\eqref{eq:pathwise-resolvent-approximation}.

Finally, the Markov semigroup representation
\[
 u_{n,\alpha}=\int_0^\infty\alpha e^{-\alpha t}
                              e^{-tG_n}g_n\dd t
\]
shows that $0\le u_{n,\alpha}\le1$. The resolvent identity
$G_nu_{n,\alpha}=\alpha(g_n-u_{n,\alpha})$ proves the remaining bound in
\eqref{eq:bounded-generator}.
\end{proof}

We now have approximations accurate for all times of the stationary path,
not merely at a fixed time. Their generator bound makes that accuracy useful
even when the starting time of an increment is selected from the path.

\subsection{Stopping-time increments}

\begin{proposition}\label{prop:path-limit}
On every fixed realization supplied by Proposition~\ref{prop:inputs}, and for
every $T>0$,
\[
 (\xi^n_t)_{0\le t\le T}\ \Longrightarrow\
 (B^h_t)_{0\le t\le T}
 \quad\hbox{in }D([0,T],\Sph)
\]
for the Skorokhod $J_1$ topology associated with the round metric $\rho$.
Every subsequential path limit is continuous.
\end{proposition}

\begin{proof}
Fix $d>0$. Choose finitely many round disks $O_1,\ldots,O_m$ covering
$\Sph$ and smooth cutoffs $f_i=1-\chi_i$ to which
Lemma~\ref{lem:uniform-recovery} applies, such that
\begin{equation}\label{eq:separating-cutoffs}
 f_i=0\text{ on }O_i,
 \qquad x\in O_i,\ \rho(x,y)\ge d\ \Longrightarrow\ f_i(y)=1.
\end{equation}
For example, take sufficiently small inner disks and slightly larger
support disks, all compactly contained in chart regions and all with
diameter smaller than $d$. Compactness gives a finite subcover, and there is
room for the recovery buffers outside each support disk.

Let $\tau,\sigma$ be stopping times for $Y^n$, with
$0\le\tau\le\sigma\le T$ and $\sigma-\tau\le s$. Let $I$ be the first index
for which $\xi^n_\tau\in O_I$. This index is measurable with respect to the
information $\mathcal F_\tau$ at time $\tau$. Since $0\le f_i\le1$,
\eqref{eq:separating-cutoffs} implies
\[
 \PP\{\rho(\xi^n_\tau,\xi^n_\sigma)\ge d\}
 \le\EE\bigl[f_I(\xi^n_\sigma)-f_I(\xi^n_\tau)\bigr].
\]
Denote by $u^{(i)}_{n,\alpha}$ the resolvent approximation of $f_i\circ z_n$.
Optional sampling of its generator martingale gives
\begin{align*}
 &\EE\bigl[u^{(I)}_{n,\alpha}(Y^n_\sigma)
                -u^{(I)}_{n,\alpha}(Y^n_\tau)\bigr]\\
 &\qquad=-\sum_{i=1}^m\EE\left[
    \mathbf1_{\{I=i\}}\int_\tau^\sigma
                    G_nu^{(i)}_{n,\alpha}(Y^n_t)\dd t\right]
 \le\alpha s.
\end{align*}
The equality is valid because each event $\{I=i\}$ lies in
$\mathcal F_\tau$; no independence of $I$ and the path is required.
The last inequality uses \eqref{eq:bounded-generator} and the fact that the
events $\{I=i\}$ partition the sample space. Consequently
\begin{align}
 \PP\{\rho(\xi^n_\tau,\xi^n_\sigma)\ge d\}
 &\le\alpha s+2\sum_{i=1}^m
 \EE\sup_{0\le t\le T}
   \abs{f_i(\xi^n_t)-u^{(i)}_{n,\alpha}(Y^n_t)}.
 \label{eq:aldous-bound}
\end{align}
This bound is uniform over all the stated stopping-time pairs. First take
$\limsup_n$, then let $s\downarrow0$ with $\alpha$ fixed, and finally let
$\alpha\to\infty$. Lemma~\ref{lem:resolvent-approximation} gives
\[
 \lim_{s\downarrow0}\limsup_n
 \sup_{\substack{\tau\le\sigma\le T\\0\le\sigma-\tau\le s}}
 \PP\{\rho(\xi^n_\tau,\xi^n_\sigma)\ge d\}=0.
\]
This is the stopping-time increment condition in Aldous's tightness
criterion \cite{Aldous1978}. Compact containment is automatic on $\Sph$,
so the laws are tight in $D([0,T],\Sph)$ with its $J_1$ topology.

Every jump of $\xi^n$ joins the endpoints of an embedded map edge.
Proposition~\ref{prop:inputs} therefore makes its largest jump tend to zero.
A jump of a $J_1$ limit must be approximated by jumps of comparable size;
hence every subsequential limit is continuous. Evaluation at fixed times is
continuous at continuous paths, so Proposition~\ref{prop:operators}
identifies the finite-dimensional laws of every such limit with those of
$B^h$. These finite-dimensional laws determine the continuous path law.
Tightness and uniqueness of the subsequential limit prove the proposition.
\end{proof}
